\chapter[Formal Problem Formulation]{Formal Problem Formulation: OASG under Degraded Observation}
\label{ch:formal-problem-formulation}

\section{State, observation, action, and constraints}
Let $x_t^\star$ denote the true physical state of the plant at time $t$, let $\hat{x}_t$ denote
the observation available to the nominal policy, and let $u_t$ denote the candidate action
proposed for release. Let $\mathcal{C}(x)$ denote the set of admissible actions when the state is
$x$, or equivalently the safety predicate that separates legal from illegal actuation. ORIUS is
concerned with regimes in which the software stack continues to evaluate legality on $\hat{x}_t$
even though the physically relevant state is $x_t^\star$ or a broader observation-consistent set
around it.

This distinction matters because physical safety is defined on the plant, not on the software
representation alone. If $\hat{x}_t$ lags, omits, or distorts the relevant state variables, then
the nominal controller can continue producing actions that are legal with respect to the observed
state while being unsafe with respect to the state that is actually realized. The formal problem
of the dissertation is therefore not generic uncertainty. It is the mismatch between
observation-conditioned legality and physical legality.

\section{Observation degradation operator and fault process}
Observation degradation is modeled as a runtime operator
\[
\hat{x}_t = \mathcal{D}(x_t^\star;\delta_t),
\]
where $\delta_t$ is a fault process that summarizes the degradation mechanisms affecting the
telemetry channel. In the repository and benchmark, this process includes stale packets, missing
samples, delay, packet reordering, spikes, drift, transport corruption, and domain-specific
state-estimation loss. The universal object of the dissertation is therefore not a plant alone;
it is the coupled pair $(x_t^\star,\hat{x}_t)$ together with the degradation process that can make
those two states diverge.

Two implications follow. First, the fault process is not merely a nuisance variable to be
averaged away in nominal performance statistics. It is part of the safety problem definition.
Second, runtime safety cannot be reduced to point estimation of the true state. Under degraded
observation, the runtime must reason over a set of plausible physical states rather than rely on
one best guess.

\section{Fault taxonomy}
The fault process is treated generically, but the benchmark and adapters instantiate it through a
shared taxonomy:
\begin{enumerate}
  \item \textbf{Freshness faults}: stale or frozen telemetry that preserves coherence while
  erasing temporal validity.
  \item \textbf{Availability faults}: missing packets, dropped samples, and partial channel loss.
  \item \textbf{Timing faults}: delay, jitter, or reordering that breaks synchronization with the
  physical plant.
  \item \textbf{Value faults}: spikes, outliers, drift, bias, or corruption that distort the
  observed state itself.
  \item \textbf{Inference faults}: estimator or learned-model errors that turn degraded inputs
  into overconfident downstream summaries.
\end{enumerate}

This taxonomy is intentionally broad enough to travel across all ORIUS domains. A battery
controller and an aerospace controller do not share the same raw sensors, but they do share the
possibility that freshness, timing, value, or inference degradation breaks the semantics of
observed legality.

\section{Definition of the observation--action safety gap}
The central event of the dissertation is the observation--action safety gap:
\[
u_t \in \mathcal{C}(\hat{x}_t)
\qquad \text{but} \qquad
u_t \notin \mathcal{C}(x_t^\star).
\]
Equivalently, the action appears safe on the observed state but unsafe on the true physical
state. This event is what the benchmark later measures through true-state safety violation and
observed-state satisfaction. It is also what makes degraded observation a first-class runtime
problem rather than an upstream modeling nuisance.

The OASG formulation is stronger than a generic claim that ``uncertainty matters.'' Uncertainty
can be large without creating a safety violation if the action remains conservative enough. The
OASG is the specific regime where the observation channel misleads the release decision itself.
That is why ORIUS focuses on action mediation, not only on uncertainty reporting.

\section{Observation-consistent state sets}
ORIUS does not assume that the true state can be recovered exactly at runtime. Instead, it works
with an observation-consistent state set
\[
\mathcal{X}_t = \left\{x : x \text{ remains plausible given } \hat{x}_t \text{ and runtime reliability } w_t\right\},
\]
where $w_t \in [0,1]$ is the observation reliability score. As reliability degrades,
$\mathcal{X}_t$ widens. This widening is the formal bridge between degraded telemetry and safe
action. Once $\mathcal{X}_t$ expands, the admissible action set must contract:
\[
\mathcal{U}_t^{\mathrm{safe}} = \{u : u \in \mathcal{C}(x) \text{ for every } x \in \mathcal{X}_t\}.
\]

The runtime question is then precise: given a candidate action $u_t$, does it belong to the
current safe set $\mathcal{U}_t^{\mathrm{safe}}$? If not, can it be repaired by projection or
replacement? If not, which governed fallback policy must be triggered? This formulation is what
lets ORIUS connect reliability scoring, uncertainty inflation, safe-set tightening, and
certificate semantics into one runtime loop.

\section{Repair, fallback, and release}
Let $\Pi_t$ denote the runtime repair map. Given a candidate action $u_t$, ORIUS computes either
a repaired action
\[
\tilde{u}_t = \Pi_t(u_t,\mathcal{X}_t)
\]
that lies inside $\mathcal{U}_t^{\mathrm{safe}}$, or else escalates to a fallback action
$u_t^{\mathrm{fb}}$ specified by the domain governance policy. The formal point is that the
runtime layer does not merely observe illegality; it mediates release. This distinguishes ORIUS
from alarm-only monitoring and from nominal predictive systems that surface confidence but leave
the release decision unchanged.

The release object is completed by a certificate
\[
\kappa_t = (\hat{x}_t, w_t, \mathcal{X}_t, \tilde{u}_t, h_t, g_t),
\]
where $h_t$ denotes a validity horizon or half-life proxy and $g_t$ denotes governance metadata
such as lifecycle state, fallback reason, or audit hash context. Later chapters make that
certificate operational through CertOS.

\section{Dual-trajectory evaluation}
The dissertation evaluates ORIUS on dual trajectories. The \emph{observed trajectory} is the one
visible to the nominal controller and therefore to the candidate action it proposes. The
\emph{true trajectory} is the one used to determine whether that candidate or repaired action
actually violates the physical safety envelope. A benchmark that inspects only observed-state
legality cannot see the OASG directly. ORIUS-Bench therefore records the candidate action, the
repaired action, true-state violation, observed-state satisfaction, optional true and observed
margins, fallback and intervention markers, latency, and certificate validity inside one replay
schema.

This dual-trajectory formulation is the critical measurement move of the dissertation. It keeps
the benchmark aligned with the formal problem object instead of allowing nominal performance
metrics to impersonate safety evidence.

\section{DomainAdapter grounding}
The universal formulation is grounded in the repository through the \texttt{DomainAdapter}
contract. Each domain must supply:
\begin{enumerate}
  \item telemetry parsing and state semantics,
  \item uncertainty or margin construction hooks,
  \item admissible-action or repair logic,
  \item true-state violation predicates,
  \item observed-state satisfaction predicates, and
  \item certificate payload specifics.
\end{enumerate}
This boundary is what makes the universal claim scientifically credible. The dissertation does
not claim that one plant model or one controller solves every domain. It claims that every domain
can expose the same runtime safety objects and thereby enter the same benchmark and governance
discipline.

\section{Scope and assumptions}
The formal problem is intentionally narrower than a claim of general AI safety. ORIUS assumes:
\begin{itemize}
  \item an actionable physical system with explicit or computable safety constraints,
  \item an observation channel that can degrade in ways that matter to release semantics,
  \item a runtime interposition layer capable of intercepting, repairing, or replacing actions,
  \item and a governance layer capable of recording why a release remained admissible.
\end{itemize}

ORIUS does not assume perfect state recovery, universal optimal control, or equal empirical
maturity across all domains. It also does not assume that uncertainty calibration alone proves
safety. The later theorem and validation chapters are therefore written under explicit assumptions
and are tied to specific artifact surfaces rather than left as narrative implication.

\paragraph{Boundary of this chapter.}
This chapter defines the formal objects that the rest of the dissertation uses: degraded
observation, observation-consistent state sets, repaired release, and dual-trajectory evaluation.
It does not yet prove the theorem ladder or claim equal cross-domain closure; those claims remain
bounded by the assumptions and evidence surfaces introduced in the chapters that follow.
